\chapter*{Glossary}
\addcontentsline{toc}{chapter}{Glossary}

% Fixed project vocabulary. The same words are used throughout the text.

\begin{description}\setlength{\itemsep}{2pt}
\item[Record] One prompt paired with the drawing made from it. Records are not linked to people. \nTotal\ were processed and \nAnalysed\ analysed.
\item[Prompt] The words a visitor typed. Prompts are quoted exactly as typed, lower case included.
\item[Drawing] The black-and-white line drawing the installation returned for a prompt.
\item[Prompt reader] The text model that reads the prompt: a RoBERTa model fine-tuned on GoEmotions (28 categories), summed into seven families.
\item[Describer] The vision model that answers five factual questions about a drawing (BLIP-VQA). It is never asked for an emotion.
\item[Image reader] The text model that infers a family from the describer's one-sentence description (BART-MNLI, zero-shot).
\item[Describe-then-extract] The image side of the pipeline: describer, template, image reader.
\item[Direct CLIP] The comparison reader, an OpenCLIP model asked for a family directly. Called the baseline; nothing else is.
\item[Category] One of the 28 GoEmotions labels.
\item[Family] One of seven groups the categories are summed into: joy, sadness, anger, fear, surprise, disgust, neutral.
\item[Agreement] The quantity this dissertation measures: how similar the prompt reader's reading of the words is to the image reader's reading of the drawing. It describes two machine readers and says nothing about what a person would see.
\item[Valence sign] $+$ for joy; $-$ for sadness, anger, fear and disgust; $0$ for surprise and neutral. Valence-sign agreement is computed over records whose prompt family has a sign.
\item[Family agreement] The share of records on which the two readers' top families coincide.
\item[Top-3 soft overlap] The overlap of the two readers' top-3 distributions once each is normalised to sum to one.
\item[Valence--arousal similarity] One minus the normalised distance between the two readings' score-weighted centroids on the circumplex. Its floor over the family centroids is 0.412, and randomly paired records already average about \FullNeuVaNull.
\item[Random-pairing null] The mean and spread a measure takes when the image readings are re-paired with the prompts at random (\nullShuffles\ re-pairings, seed \nullSeed). Every reported value is read against it.
\item[$z$] (observed $-$ null mean) / null standard deviation. Used to order effects, not as a precise probability: across seeds it moves by about a unit (\ZSeedMin\ to \ZSeedMax\ for the headline).
\item[Degenerate null] A null with zero spread. It arises when every prompt in a subset shares one label, so re-pairing cannot change the score. No $z$ is reported for it.
\item[Prompt group] All records whose prompt text is identical after case-folding, computed before redaction. There are \nDistinct\ groups; \nRepeats\ records repeat an earlier prompt. Every interval and test clusters on the group.
\item[Repeat group] A prompt group with two or more records (\nGroups\ of them, \nGroupRecords\ records).
\item[Clustered bootstrap] Resampling prompt groups, with replacement, to get 95\% confidence intervals that respect repeated prompts (\DiffDraws\ draws).
\item[Face-classified / scene-classified] A record is face-classified when the describer reports at least one person \emph{and} its expression answer contains a readable expression word. Every other record is scene-classified. The split is the describer's, not a person's.
\item[Emotional / neutral-labelled prompt] A prompt is emotional when the prompt reader's top family is not neutral (\nEmotional\ records); otherwise it is neutral-labelled (\nNonEmotional). The label was never checked by a person.
\item[Specificity check] The expectation that neutral-labelled prompts show no agreement above their null. Only valence--arousal similarity can move there.
\item[Copy rate] Where the describer's expression answer is itself an emotion word, the share of records in which the image reader's top family is that word's family (\DescCopyStrict\ on the strict definition).
\item[Consistency] For repeat groups, the mean share of drawings in a group that the image reader gives the group's most common family, read against a permutation null.
\item[People-aware template] The ablation variant of the description template that leaves out the person and expression clauses wherever the describer reported no people.
\item[Single-coder check] The 120-drawing sample the author codes blind, from contact sheets, to score the describer and the image reader against a person.
\item[Demonstrator] The local web application that runs the real pipeline on any record or an uploaded image and shows every measure against its chance level.
\end{description}
